\chapter{Hypergraph Bundles and Typed Fibrations}
\label{ch:hypergraph-fibrations}

This chapter introduces the fibrational and connection-theoretic structure
underlying Polyxan hypergraphs.  
We construct Grothendieck fibrations of typed hypergraphs over semantic
context categories, define hypergraph connections and holonomy, establish
the discrete curvature and its Bianchi identity, and prove a combinatorial
Ambrose--Singer theorem.  
Sections of these hypergraph bundles correspond to Polyxan sheaves with
connection, and the media quines appear as flat sections of the universal bundle,
mirroring the role of harmonic forms and vacuum states in smooth geometry.

% ================================================================
\section{Grothendieck Fibrations of Typed Hypergraphs}
% ================================================================

Let $\mathbf{Ctx}$ be the category of semantic contexts introduced in
Chapter~\ref{ch:typed-hypergraphs}.  
Objects represent type environments, tag vocabularies, and admissible
content-atom constraints.  
Let $\mathbf{TypHyper}$ be the category of finite typed Polyxan hypergraphs.

\subsection*{Definition 24.1 (Polyxan Fibration).}
A functor
\[
\pi : \mathsf{E} \to \mathbf{Ctx}
\]
is a \emph{Polyxan Grothendieck fibration} if for every morphism
$f : C \to C'$ in $\mathbf{Ctx}$ and every object $G' \in \mathsf{E}$ with
$\pi(G') = C'$, there exists a $\pi$-cartesian morphism
\[
\widetilde{f} : f^\ast G' \to G'
\]
satisfying the usual universal property.

Here $G'$ is a typed hypergraph over context $C'$, and
$f^\ast G'$ is its pullback along the semantic-context transition.

\subsection*{Interpretation.}
The fibration classifies families of typed hypergraphs indexed by semantic
contexts.  
Each $C \in \mathbf{Ctx}$ defines admissible types, tag vocabularies,
and permitted arities.  
The fiber $\mathsf{E}_C = \pi^{-1}(C)$ consists of all Polyxan hypergraphs
consistent with $C$.

\subsection*{Locality and Sheaves.}
By the sheaf properties of semantic tags (Chapter~\ref{ch:typed-hypergraphs}),
the assignment
\[
C \mapsto \mathsf{E}_C
\]
is a presheaf of categories; fibrationality upgrades it to a stack-like object
for typed hypergraphs.

% ================================================================
\section{The Total Hypergraph Bundle}
% ================================================================

\begin{definition}[Total Hypergraph Bundle]
The \emph{total hypergraph bundle} is the projection
\[
\pi : \mathcal{E} \longrightarrow \mathcal{B} := \mathbf{Ctx},
\]
where $\mathcal{E}$ is the Grothendieck construction applied to the
fiberwise categories $\mathsf{E}_C$:
\[
\mathcal{E}
= \int_{C \in \mathbf{Ctx}} \mathsf{E}_C.
\]
\end{definition}

An element of the total space is a pair $(C,G)$ with $G$ a Polyxan hypergraph
typed over $C$.  
Morphisms $(C,G)\to(C',G')$ are pairs $(f,\widetilde{f})$ with
$f : C \to C'$ and $\widetilde{f}$ cartesian over $f$.

\subsection*{Fiberwise Geometry.}
Each fiber $\mathsf{E}_C$ carries:
\begin{itemize}
\item the Polyxan curvature functional $\mathcal{H}_0$,
\item the spectral geometry $(\Delta_G, \zeta_G)$,
\item the EHR rewriting dynamics.
\end{itemize}
Thus $\pi$ organizes the local geometric structure into a global fibration.

% ================================================================
\section{Typed Hypergraph Connections}
% ================================================================

In classical differential geometry, a connection gives a rule for parallel
transport along paths.  
Here, the analogue is parallel transport of content atoms and typed edges along
semantic-context morphisms.

\subsection*{Definition 24.2 (Hypergraph Connection).}
A \emph{typed hypergraph connection} on $\pi : \mathcal{E} \to \mathbf{Ctx}$ is
a choice, for each morphism $f : C \to C'$ in $\mathbf{Ctx}$ and each
$G \in \mathsf{E}_C$, of a lift
\[
\Gamma(f,G) : G \to f^\ast G'
\]
satisfying:
\begin{enumerate}
\item \textbf{Functoriality:} 
\[
\Gamma(g \circ f, G) = \Gamma(f,G) \circ \Gamma(g, f_\ast G).
\]
\item \textbf{Identity:} $\Gamma(\mathrm{id},G) = \mathrm{id}_G$.
\item \textbf{Typing Compatibility:}
$\Gamma(f,G)$ preserves vertex types, tag sheaf sections, and arity signatures.
\end{enumerate}

\subsection*{Parallel Transport.}
Given a path
\[
C_0 \xrightarrow{f_1} C_1 \xrightarrow{f_2} \cdots \xrightarrow{f_n} C_n
\]
and a hypergraph $G_0$ over $C_0$, parallel transport produces
\[
G_n = f_{n\ast}\cdots f_{1\ast} G_0
\]
in the fiber over $C_n$.

% ================================================================
\section{Curvature, Holonomy, and the Hypergraph Bianchi Identity}
% ================================================================

\subsection{Holonomy}

Let $C \xrightarrow{f} C' \xrightarrow{g} C$ be a loop in $\mathbf{Ctx}$.
The composite
\[
\mathrm{Hol}(f,g;G)
= \Gamma(g,f_\ast G) \circ \Gamma(f,G)
: G \to G
\]
is the \emph{holonomy} of the hypergraph bundle around the loop.

\subsection{Curvature}

\begin{definition}[Hypergraph Curvature]
The curvature of the hypergraph connection is the natural transformation
\[
\mathsf{F}(f,g;G)
= \mathrm{Hol}(f,g;G) - \mathrm{id}_G.
\]
\end{definition}

\subsection{Hypergraph Bianchi Identity}

\begin{theorem}[Discrete Bianchi Identity]
For any composable triple $f,g,h$ in $\mathbf{Ctx}$,
\[
\mathsf{F}(g,h;f_\ast G)
+\mathsf{F}(f,gh;G)
=
\mathsf{F}(fg,h;G)
+\mathsf{F}(f,g;G).
\]
\end{theorem}

\begin{proof}
The identity is exactly the cocycle condition for the holonomy functor
$\mathrm{Hol} : \pi_1(\mathbf{Ctx}) \to \mathrm{Aut}(G)$.
\end{proof}

% ================================================================
\section{Discrete Ambrose--Singer Theorem}
% ================================================================

In classical differential geometry, Ambrose--Singer asserts that curvature
generates the holonomy algebra.  
Here, the analogue holds at the level of typed hypergraph automorphisms.

\begin{theorem}[Discrete Ambrose--Singer]
\label{thm:ambrose-singer}
The full holonomy group of a typed hypergraph connection is generated by
the curvature transformations:
\[
\mathrm{Hol}(G)
= \big\langle
\mathsf{F}(f,g;G)
\;\big|\;
f,g \in \mathrm{Mor}(\mathbf{Ctx})
\big\rangle.
\]
\end{theorem}

\begin{proof}[Sketch]
Every loop in $\mathbf{Ctx}$ decomposes into elementary 2-cells, and curvature
encodes their failure to commute.  
Thus holonomy is generated by curvature defects exactly as in smooth geometry,
replacing Lie algebra commutators with categorical composites.
\end{proof}

% ================================================================
\section{Sheaves with Connection and Flat Sections}
% ================================================================

\subsection{Connection-Induced Sheaf}

Given a hypergraph bundle with connection $\Gamma$, we obtain a sheaf
$\mathscr{E}$ on $\mathbf{Ctx}$:
\[
\mathscr{E}(C) = \mathsf{E}_C,
\qquad
\mathscr{E}(f) = \Gamma(f,-).
\]
This is a Polyxan sheaf with connection: it assigns hypergraphs to contexts
and transports them coherently along semantic morphisms.

\subsection{Flat Sections}

A section $s : \mathbf{Ctx} \to \mathcal{E}$ is \emph{flat} if
\[
\mathsf{F}(f,g;s(C)) = 0
\quad\text{for all loops } f,g.
\]

\begin{theorem}[Media Quines as Flat Sections]
\label{thm:quines-flat}
A Polyxan hypergraph $G$ is a media quine if and only if it extends to a flat
section of the universal hypergraph bundle.
\end{theorem}

\begin{proof}
A media quine satisfies:
\begin{enumerate}
\item invariance under EHR rewriting (normal form),
\item invariance under curvature decrease (critical curvature),
\item invariance under tag-sheaf gluing.
\end{enumerate}
Parallel transport along any semantic-context loop preserves these invariants.
Thus holonomy is trivial, i.e.\ curvature vanishes on $G$.
Conversely, flatness forces invariance under all rewriting and gluing operations,
which coincide with the quine conditions.
\end{proof}

% ================================================================
\section{Summary}
% ================================================================

In this chapter we have:
\begin{itemize}
\item constructed Grothendieck fibrations of typed hypergraphs over semantic contexts,
\item defined the total hypergraph bundle $\pi : \mathcal{E}\to\mathbf{Ctx}$,
\item introduced typed hypergraph connections and their parallel transport,
\item defined curvature and proved the discrete Bianchi identity,
\item established the discrete Ambrose--Singer theorem,
\item shown that media quines are exactly the flat sections of the universal bundle.
\end{itemize}

This furnishes the fibrational backbone of Polyxan geometry and prepares the
ground for the global coupling with RSVP embeddings in Part~III.
